% Source: https://artofproblemsolving.com/community/c3h1482056_sliding_triangles
% Stable identity: aops-topic-1482056-q01
An isosceles triangle with two sides of length $20$ and base $24$ lies on a flat surface (see the figure). Its top vertex is pushed down, so that the two equal sides keep their length $20$, the angle between them increases and the base stretches, until the triangle collapses. During this process the area first increases and then decreases.

\begin{asy}
usepackage("amsmath");
size(4.5cm);
defaultpen(fontsize(10pt));
pair A=(-12,0),B=(12,0),C=(0,16);
draw(A--B--C--cycle,linewidth(0.9));
dot(A);dot(B);dot(C);
label("$20$",(A+C)/2,NW); label("$20$",(B+C)/2,NE); label("$24$",(A+B)/2,S);
label("\mbox{}",(min(currentpicture,user=true)+max(currentpicture,user=true))/2,black); // keeps the text after the figure black when LaTeX typesets the labels inline
\end{asy}

\begin{enumerate}
\item[(a)] What is the maximum area reached during the process?
\item[(b)] At one other moment the triangle has the same area as the original triangle. What is the length of the base at that moment?
\end{enumerate}
